\chapter*{부 ~ ~ ~록}\label{chap:appendix}
\addcontentsline{toc}{chapterstar}{부록}
\vspace*{4em}
\anset

학생들을 대상으로 투입한 문항지는 다음과 같다. \\

%%%%%%%%%%%%%%%%%%%%%%%%%%%%%%%%%%%%%%%%%%%%%%%%%%%%%%%%%%%%%%%%%%%%%%%%%%%%

\bnset
\begin{graybox}\centering
{\fontsize{14pt}{17pt}\selectfont\bd{\an
\ItemTitlei
}}\vspace*{0.8em}\end{graybox}\vspace{0.4em}

\begin{graybox}\vspace*{0.2em}\bn
눈금 간격이 0.1\,cm인 자를 사용했을 때, 측정값은 소수점 아래 몇 자리까지 쓰는 것이
적절하다고 생각하나요?
(단답형) \\ 답: 
\vspace*{0.4em}\end{graybox}\vspace{0.4em}

\begin{graybox}\vspace*{0.2em}\bn
위에서 그렇게 답변한 이유를 말해주세요.
(서술형, 30자 이상) \\ 서술: 
\vspace*{2.4em}\end{graybox}\vspace{0.4em}

%%%%%%%%%%%%%%%%%%%%%%%%%%%%%%%%%%%%%%%%%%%%%%%%%%%%%%%%%%%%%%%%%%%%%%%%%%%%

\clearpage

%%%%%%%%%%%%%%%%%%%%%%%%%%%%%%%%%%%%%%%%%%%%%%%%%%%%%%%%%%%%%%%%%%%%%%%%%%%%

\bnset
\begin{graybox}\centering
{\fontsize{14pt}{17pt}\selectfont\bd{\an
\ItemTitleii
}}\vspace*{0.8em}\end{graybox}\vspace{0.4em}

\begin{graybox}
~~ 다음은 자로 연필의 길이를 측정한 모습을 나타낸 것입니다. 연필의 왼쪽 끝과 오른쪽 끝이
자의 몇 cm 눈금에 있는지 각각 쓰고, 두 눈금의 차이를 구해 연필의 길이를 쓰세요.
\rulerpencilnotag
\end{graybox}\vspace{0.4em}

\begin{graybox}\vspace*{0.2em}\bn
연필 왼쪽 끝 눈금의 위치를 쓰세요.
(단답형) \\ 답: 
\vspace*{0.4em}\end{graybox}\vspace{0.4em}

\begin{graybox}\vspace*{0.2em}\bn
연필 오른쪽 끝 눈금의 위치를 쓰세요.
(단답형) \\ 답: 
\vspace*{0.4em}\end{graybox}\vspace{0.4em}

\begin{graybox}\vspace*{0.2em}\bn
연필의 길이를 쓰세요.
(단답형) \\ 답: 
\vspace*{0.4em}\end{graybox}\vspace{0.4em}

\begin{graybox}\vspace*{0.2em}\bn
연필의 오른쪽 끝은 눈금 사이 어디쯤 걸쳐 있습니다.
이런 경우, 측정값을 소수점 아래 몇 자리까지 적는 것이
적절하다고 생각하는지 작성해주세요.
(단답형) \\ 답: 
\vspace*{0.4em}\end{graybox}\vspace{0.4em}

\begin{graybox}\vspace*{0.2em}\bn
위에서 그렇게 답변한 이유를 말해주세요.
(서술형, 30자 이상) \\ 서술: 
\vspace*{2.4em}\end{graybox}\vspace{0.4em}

%%%%%%%%%%%%%%%%%%%%%%%%%%%%%%%%%%%%%%%%%%%%%%%%%%%%%%%%%%%%%%%%%%%%%%%%%%%%

\clearpage

%%%%%%%%%%%%%%%%%%%%%%%%%%%%%%%%%%%%%%%%%%%%%%%%%%%%%%%%%%%%%%%%%%%%%%%%%%%%

\bnset
\begin{graybox}\centering
{\fontsize{14pt}{17pt}\selectfont\bd{\an
\ItemTitleiii
}}\vspace*{0.8em}\end{graybox}\vspace{0.4em}

\begin{graybox}
~~ 다음은 버니어캘리퍼스를 이용하여 연필의 길이를 측정하는 모습을 나타낸 것입니다.
두 학생이 연필의 길이를 각각 측정한 뒤, 측정 결과를 다음과 같이 말하였습니다.
\begin{center}학생 A : 8.628\,cm \\ 학생 B : 8.625\,cm\end{center}
\caliperpencilnotg
\end{graybox}\vspace{0.4em}

\begin{graybox}\vspace*{0.2em}\bn
두 학생이 측정한 연필 길이의 평균을 구해보세요.
(단답형) \\ 답: 
\vspace*{0.4em}\end{graybox}\vspace{0.4em}

\begin{graybox}\vspace*{0.2em}\bn
평균값을 기록할 때 소수점 아래 몇 자리까지 쓰는 것이 적절하다고 생각하나요?
사용하는 도구의 정밀도를 고려하여 판단해 보세요.
(단답형) \\ 답: 
\vspace*{0.4em}\end{graybox}\vspace{0.4em}

%%%%%%%%%%%%%%%%%%%%%%%%%%%%%%%%%%%%%%%%%%%%%%%%%%%%%%%%%%%%%%%%%%%%%%%%%%%%

\clearpage

%%%%%%%%%%%%%%%%%%%%%%%%%%%%%%%%%%%%%%%%%%%%%%%%%%%%%%%%%%%%%%%%%%%%%%%%%%%%

\bnset
\begin{graybox}\centering
{\fontsize{14pt}{17pt}\selectfont\bd{\an
\ItemTitleiv
}}\vspace*{0.8em}\end{graybox}\vspace{0.4em}

\begin{graybox}
~~ 다음은 9명의 학생이 자를 이용하여 연필의 길이를 측정한 결과를 나타낸 것입니다.
9개의 값을 계산기로 평균한 값은 8.64333333....\,cm 입니다.
\SeveralMeasurementValues
\end{graybox}\vspace{0.4em}

\begin{graybox}\vspace*{0.2em}\bn
위의 측정값 평균은 소수점 아래 몇 자리까지 쓰는 것이 적절하다고 생각하나요?
(단답형) \\ 답: 
\vspace*{0.4em}\end{graybox}\vspace{0.4em}

\begin{graybox}\vspace*{0.2em}\bn
선택한 자릿수의 근거를 도구의 정밀도와 연결지어 설명해 보세요.
(서술형, 30자 이상) \\ 서술: 
\vspace*{2.4em}\end{graybox}\vspace{0.4em}

%%%%%%%%%%%%%%%%%%%%%%%%%%%%%%%%%%%%%%%%%%%%%%%%%%%%%%%%%%%%%%%%%%%%%%%%%%%%

\clearpage

%%%%%%%%%%%%%%%%%%%%%%%%%%%%%%%%%%%%%%%%%%%%%%%%%%%%%%%%%%%%%%%%%%%%%%%%%%%%

\bnset
\begin{graybox}\centering
{\fontsize{14pt}{17pt}\selectfont\bd{\an
\ItemTitlev
}}\vspace*{0.8em}\end{graybox}\vspace{0.4em}

\begin{graybox}
~~ 다음은 쿨롱 상수 $k$를 얻는 과정입니다. \vspace*{10pt}
\begin{graybox}\vspace*{0.6em}
\bul 진공 중에서 빛의 속도 $c$\,는 대략 $299\,\,792\,\,458\,\mathrm{m/s}$이다.
\bul 진공 중의 유전율 $\epsilon_0 = \dfrac{1}{c^2 \mu_0}$은 ... 대략
$8.854\,\,187\,\,818\,\,8 \times 10^{-12}\,\mathrm{F \cdot m^{-1}}$이다.
\bul 쿨롱 상수 $k = \dfrac{1}{4\pi \epsilon_0}$은 ... 대략
$8.987\,\,551\,\,786\,\,2 \times 10^9\,\mathrm{N \cdot m^2/C^2}$으로 나타낼 수 있다.
\vspace*{0.6em}\end{graybox}
\end{graybox}\vspace{0.4em}

\begin{graybox}\vspace*{0.2em}\bn
거리가 1.5\,m 떨어진 두 전하의 전하량이 각각 2.0\,C,\, 4.5\,C일 때,
쿨롱의 법칙을 이용하여 두 전하에 작용하는 힘을 구하고자 합니다.
계산 과정에서 쿨롱 상수를 사용할 때 소수점 아래 몇 자리까지 사용하는 것이 적절하다고 생각하나요?
계산 결과의 정밀도를 고려하여 자릿수를 판단하고 이유를 설명해 보세요.
(단답형) \\ 답: 
\vspace*{0.4em}\end{graybox}\vspace{0.4em}

\begin{graybox}\vspace*{0.2em}\bn
위에서 그렇게 답변한 이유를 말해주세요.
(서술형, 30자 이상) \\ 서술: 
\vspace*{2.4em}\end{graybox}\vspace{0.4em}

%%%%%%%%%%%%%%%%%%%%%%%%%%%%%%%%%%%%%%%%%%%%%%%%%%%%%%%%%%%%%%%%%%%%%%%%%%%%

\clearpage

%%%%%%%%%%%%%%%%%%%%%%%%%%%%%%%%%%%%%%%%%%%%%%%%%%%%%%%%%%%%%%%%%%%%%%%%%%%%

\begin{graybox}
~~ 계산에서 쿨롱 상수를 정해진 자릿수보다 더 많거나 적은 자릿수로 사용할 경우, 결과값에
어떤 과학적 영향이 생길 수 있는지 설명해 보세요. 오차, 정밀도, 신뢰도 측면에서 판단해
보세요.
\end{graybox}\vspace{0.4em}

\begin{graybox}\vspace*{0.2em}\bn
더 많은 자릿수를 사용할 경우 결과에 영향이 있다고 생각하는지 선택해주세요.
(선택형) \\
\raisebox{-0.25ex}{\scalebox{1.4}{$\Box$}} 더 많은 자릿수를 선택하면 결과에 영향이 있다. \\
\raisebox{-0.25ex}{\scalebox{1.4}{$\Box$}} 더 많은 자릿수를 선택하면 결과에 영향이 없다.
\vspace*{0.4em}\end{graybox}\vspace{0.4em}

\begin{graybox}\vspace*{0.2em}\bn
위에서 그렇게 답변한 이유를 말해주세요.
(서술형, 30자 이상) \\ 서술: 
\vspace*{2.4em}\end{graybox}\vspace{0.4em}

\begin{graybox}\vspace*{0.2em}\bn
더 적은 자릿수를 사용할 경우 결과에 영향이 있다고 생각하는지 선택해주세요.
(선택형) \\
\raisebox{-0.25ex}{\scalebox{1.4}{$\Box$}} 더 많은 자릿수를 선택하면 결과에 영향이 있다. \\
\raisebox{-0.25ex}{\scalebox{1.4}{$\Box$}} 더 많은 자릿수를 선택하면 결과에 영향이 없다.
\vspace*{0.4em}\end{graybox}\vspace{0.4em}

\begin{graybox}\vspace*{0.2em}\bn
위에서 그렇게 답변한 이유를 말해주세요.
(서술형, 30자 이상) \\ 서술: 
\vspace*{2.4em}\end{graybox}\vspace{0.4em}

%%%%%%%%%%%%%%%%%%%%%%%%%%%%%%%%%%%%%%%%%%%%%%%%%%%%%%%%%%%%%%%%%%%%%%%%%%%%

\clearpage

%%%%%%%%%%%%%%%%%%%%%%%%%%%%%%%%%%%%%%%%%%%%%%%%%%%%%%%%%%%%%%%%%%%%%%%%%%%%

\bnset
\begin{graybox}\centering
{\fontsize{14pt}{17pt}\selectfont\bd{\an
\ItemTitlevi
}}\vspace*{0.8em}\end{graybox}\vspace{0.4em}

\begin{graybox}\vspace*{0.2em}\bn
자로 연필의 길이를 측정하는 상황에서 유효숫자를 반드시 지켜야 한다고 생각하나요?
(선택형) \\
\raisebox{-0.25ex}{\scalebox{1.4}{$\Box$}} 유효숫자를 지켜야 한다. \\
\raisebox{-0.25ex}{\scalebox{1.4}{$\Box$}} 유효숫자를 지킬 필요 없다.
\vspace*{0.4em}\end{graybox}\vspace{0.4em}

\begin{graybox}\vspace*{0.2em}\bn
위에서 그렇게 답변한 이유를 말해주세요.
(서술형, 30자 이상) \\ 서술: 
\vspace*{2.4em}\end{graybox}\vspace{0.4em}

\begin{graybox}\vspace*{0.2em}\bn
버니어캘리퍼스로 연필의 길이를 측정하는 상황에서 유효숫자를 반드시 지켜야 한다고 생각하나요?
(선택형) \\
\raisebox{-0.25ex}{\scalebox{1.4}{$\Box$}} 유효숫자를 지켜야 한다. \\
\raisebox{-0.25ex}{\scalebox{1.4}{$\Box$}} 유효숫자를 지킬 필요 없다.
\vspace*{0.4em}\end{graybox}\vspace{0.4em}

\begin{graybox}\vspace*{0.2em}\bn
위에서 그렇게 답변한 이유를 말해주세요.
(서술형, 30자 이상) \\ 서술: 
\vspace*{2.4em}\end{graybox}\vspace{0.4em}

\begin{graybox}\vspace*{0.2em}\bn
물리 상수를 이용하여 계산할 때 유효숫자를 반드시 지켜야 한다고 생각하나요?
(선택형) \\
\raisebox{-0.25ex}{\scalebox{1.4}{$\Box$}} 유효숫자를 지켜야 한다. \\
\raisebox{-0.25ex}{\scalebox{1.4}{$\Box$}} 유효숫자를 지킬 필요 없다.
\vspace*{0.4em}\end{graybox}\vspace{0.4em}

\begin{graybox}\vspace*{0.2em}\bn
위에서 그렇게 답변한 이유를 말해주세요.
(서술형, 30자 이상) \\ 서술: 
\vspace*{2.4em}\end{graybox}\vspace{0.4em}